% !TeX program = pdflatex
\documentclass[11pt,a4paper]{amsart}

\newcommand{\notelanguage}{finnish}
% Shared preamble for course notes and exercise sets.

\usepackage[T1]{fontenc}
\usepackage{lmodern}
\providecommand{\notelanguage}{english}
\usepackage[finnish,english]{babel}
\AtBeginDocument{\expandafter\selectlanguage\expandafter{\notelanguage}}

\newcommand{\notesfinnishlanguage}{finnish}
\ifx\notelanguage\notesfinnishlanguage
\newcommand{\notestheoremname}{Lause}
\newcommand{\noteslemmaname}{Lemma}
\newcommand{\notespropositionname}{Propositio}
\newcommand{\notescorollaryname}{Korollaari}
\newcommand{\notesdefinitionname}{Määritelmä}
\newcommand{\notesexamplename}{Esimerkki}
\newcommand{\notesproblemname}{Tehtävä}
\newcommand{\notesexercisename}{Tehtävä}
\newcommand{\notesremarkname}{Huomautus}
\newcommand{\notessolutionname}{Ratkaisu}
\else
\newcommand{\notestheoremname}{Theorem}
\newcommand{\noteslemmaname}{Lemma}
\newcommand{\notespropositionname}{Proposition}
\newcommand{\notescorollaryname}{Corollary}
\newcommand{\notesdefinitionname}{Definition}
\newcommand{\notesexamplename}{Example}
\newcommand{\notesproblemname}{Problem}
\newcommand{\notesexercisename}{Exercise}
\newcommand{\notesremarkname}{Remark}
\newcommand{\notessolutionname}{Solution}
\fi

\usepackage[a4paper,margin=30mm]{geometry}
\usepackage{microtype}
\usepackage{graphicx}
\usepackage{enumitem}

\newcommand{\exerciseheading}[2]{%
  \par\noindent
  {\large\bfseries #1\par}
  \vspace{0.1em}
  \noindent{\large #2\par}
  \vspace{1.5em}
}

\usepackage{amsmath}
\usepackage{amssymb}
\usepackage{amsthm}
\usepackage{mathtools}
\usepackage{aliascnt}

% Row operations beside the affected row, e.g. \rowop{+2R_1}.
\newcommand{\rowop}[1]{\mathrel{\big|}\,#1}

% Linear systems: the argument gives the number of alignment pairs.
\newenvironment{linearsystem}[1]
{\left\{
  \begin{alignedat}{#1}}
    {
  \end{alignedat}\right.}

  \usepackage[hidelinks,pdfusetitle]{hyperref}

  \numberwithin{equation}{section}
  \setlist{itemsep=0.25em,topsep=0.5em}

  \theoremstyle{plain}
  \newtheorem{theorem}{\notestheoremname}[section]
  \newaliascnt{lemma}{theorem}
  \newtheorem{lemma}[lemma]{\noteslemmaname}
  \aliascntresetthe{lemma}
  \newaliascnt{proposition}{theorem}
  \newtheorem{proposition}[proposition]{\notespropositionname}
  \aliascntresetthe{proposition}
  \newaliascnt{corollary}{theorem}
  \newtheorem{corollary}[corollary]{\notescorollaryname}
  \aliascntresetthe{corollary}

  \theoremstyle{definition}
  \newaliascnt{definition}{theorem}
  \newtheorem{definition}[definition]{\notesdefinitionname}
  \aliascntresetthe{definition}
  \newaliascnt{example}{theorem}
  \newtheorem{example}[example]{\notesexamplename}
  \aliascntresetthe{example}
  \newaliascnt{problem}{theorem}
  \newtheorem{problem}[problem]{\notesproblemname}
  \aliascntresetthe{problem}
  \newtheorem{exercise}{\notesexercisename}

  \theoremstyle{remark}
  \newaliascnt{remark}{theorem}
  \newtheorem{remark}[remark]{\notesremarkname}
  \aliascntresetthe{remark}

  \newenvironment{solution}
  {
  \begin{proof}[\notessolutionname]}
    {
  \end{proof}}

  \usepackage[nameinlink,noabbrev]{cleveref}

  \crefname{theorem}{theorem}{theorems}
  \crefname{lemma}{lemma}{lemmas}
  \crefname{proposition}{proposition}{propositions}
  \crefname{corollary}{corollary}{corollaries}
  \crefname{definition}{definition}{definitions}
  \crefname{example}{example}{examples}
  \crefname{problem}{problem}{problems}
  \crefname{exercise}{exercise}{exercises}
  \crefname{remark}{remark}{remarks}

  \ifx\notelanguage\notesfinnishlanguage
  \crefname{theorem}{lause}{lauseet}
  \crefname{lemma}{lemma}{lemmat}
  \crefname{proposition}{propositio}{propositiot}
  \crefname{corollary}{korollaari}{korollaarit}
  \crefname{definition}{määritelmä}{määritelmät}
  \crefname{example}{esimerkki}{esimerkit}
  \crefname{problem}{tehtävä}{tehtävät}
  \crefname{exercise}{tehtävä}{tehtävät}
  \crefname{remark}{huomautus}{huomautukset}
  \fi

  \newcommand{\NN}{\mathbb{N}}
  \newcommand{\NNone}{\mathbb{N}_1}
  \newcommand{\ZZ}{\mathbb{Z}}
  \newcommand{\QQ}{\mathbb{Q}}
  \newcommand{\RR}{\mathbb{R}}
  \newcommand{\CC}{\mathbb{C}}
  \newcommand{\PP}{\mathbb{P}}

  \DeclareMathOperator{\im}{im}
  \DeclareMathOperator{\rank}{rank}
  \DeclareMathOperator{\syt}{syt}
  \DeclareMathOperator{\pyj}{pyj}
  \DeclareMathOperator{\dom}{dom}
  \DeclarePairedDelimiter{\abs}{\lvert}{\rvert}
  \DeclarePairedDelimiter{\norm}{\lVert}{\rVert}
  \newcommand{\set}[1]{\left\lbrace#1\right\rbrace}
  \newcommand{\maxset}[1]{\max\set{#1}}
  \newcommand{\minset}[1]{\min\set{#1}}
  \newcommand{\powerset}[1]{\mathcal{P}\!\left(#1\right)}


\title{Johdatus logiikkaan II\\Harjoituksen 3 ratkaisut}
\author{}
\date{}

\begin{document}

\exerciseheading{Johdatus logiikkaan II}{Harjoituksen 3 ratkaisut}

\begin{exercise}
  Osoita, että kaavat \(\forall x\exists y(R(x,y) \land R(y,x))\) ja
  \(\forall x(\exists yR(x,y) \land \exists yR(y,x))\) eivät ole loogisesti
  ekvivalentteja. Onko jompi kumpi toisen looginen seuraus?
\end{exercise}

\begin{solution}
  Merkitään
  \[
    \begin{aligned}
      A &= \forall x\exists y(R(x,y) \land R(y,x)), \\
      B &= \forall x(\exists yR(x,y) \land \exists yR(y,x)).
    \end{aligned}
  \]
  Osoitetaan, että \(B\) on kaavan \(A\) looginen seuraus.
  Olkoon \(\mathcal M\) mielivaltainen malli ja \(s\) sen tulkintafunktio.
  Oletetaan, että \(\mathcal M \models_s A\).

  Olkoon \(b \in \dom(\mathcal M)\) mielivaltainen. Universaali- ja
  eksistenssikvanttorin totuusmääritelmien nojalla on olemassa
  \(a \in \dom(\mathcal M)\), jolla
  \[
    \mathcal M \models_{s(b/x)(a/y)}
    \bigl(R(x,y) \land R(y,x)\bigr).
  \]
  Tässä alkio \(a\) saa riippua valitusta alkiosta \(b\).
  Konjunktion totuusmääritelmän nojalla
  \[
    \mathcal M \models_{s(b/x)(a/y)} R(x,y)
    \quad\text{ja}\quad
    \mathcal M \models_{s(b/x)(a/y)} R(y,x).
  \]
  Sama alkio \(a\) kelpaa siis todistajaksi kumpaankin
  eksistenssikvantifioituun kaavaan, joten
  \[
    \mathcal M \models_{s(b/x)} \exists yR(x,y)
    \quad\text{ja}\quad
    \mathcal M \models_{s(b/x)} \exists yR(y,x).
  \]
  Näin ollen
  \[
    \mathcal M \models_{s(b/x)}
    \bigl(\exists yR(x,y) \land \exists yR(y,x)\bigr).
  \]
  Koska \(b\) oli mielivaltainen, universaalikvanttorin totuusmääritelmästä
  seuraa \(\mathcal M \models_s B\). Mallin \(\mathcal M\) ja
  tulkintafunktion \(s\) mielivaltaisuuden nojalla siis \(A \models B\).

  Käänteinen seuraussuhde ei päde. Tarkastellaan mallia \(\mathcal M\), jossa
  \[
    \dom(\mathcal M) = \{1,2,3\},
    \qquad
    R^{\mathcal M} = \{(1,3),(3,2),(2,1)\}.
  \]
  Olkoon \(s\) mielivaltainen \(\mathcal M\)-tulkintafunktio.
  Jokaiselle alkiolle \(b \in \dom(\mathcal M)\) voidaan valita alkiot
  \(a,c \in \dom(\mathcal M)\) seuraavasti:
  \[
    \begin{array}{c|cc}
      b & a & c \\
      \hline
      1 & 3 & 2 \\
      2 & 1 & 3 \\
      3 & 2 & 1
    \end{array}
  \]
  Tällöin \((b,a) \in R^{\mathcal M}\) ja \((c,b) \in R^{\mathcal M}\),
  joten
  \[
    \mathcal M \models_{s(b/x)}
    \bigl(\exists yR(x,y) \land \exists yR(y,x)\bigr).
  \]
  Koska tämä pätee jokaiselle \(b\), saadaan \(\mathcal M \models_s B\).

  Kaava \(A\) ei sen sijaan ole tosi tässä mallissa. Kun \(x\):n arvoksi
  valitaan \(1\), ainoa mahdollinen todistaja kaavalle \(R(x,y)\) on
  \(y = 3\). Koska \((3,1) \notin R^{\mathcal M}\), tämä alkio ei kuitenkaan
  toteuta kaavaa \(R(y,x)\). Siten mikään \(y\):n arvo ei toteuta molempia
  kaavoja yhtä aikaa, eli
  \[
    \mathcal M \not\models_{s(1/x)}
    \exists y\bigl(R(x,y) \land R(y,x)\bigr).
  \]
  Siis \(\mathcal M \not\models_s A\), joten \(B \not\models A\).
  Näin ollen kaavat \(A\) ja \(B\) eivät ole loogisesti ekvivalentteja:
  \(B\) on kaavan \(A\) looginen seuraus, mutta \(A\) ei ole kaavan \(B\)
  looginen seuraus.
\end{solution}

\newpage

\begin{exercise}
  Osoita, että kaava \(\forall y\exists xA\) on kaavan
  \(\exists x\forall yA\) looginen seuraus (huomaa, että ei ole väliä,
  mikä kaava \(A\) on).
\end{exercise}

\begin{solution}
  Olkoon \(\mathcal M\) mielivaltainen malli ja \(s\) sen tulkintafunktio.
  Oletetaan, että
  \[
    \mathcal M \models_s \exists x\forall yA.
  \]
  Eksistenssi- ja universaalikvanttorin totuusmääritelmien nojalla on
  olemassa \(b \in \dom(\mathcal M)\), jolla kaikilla
  \(a \in \dom(\mathcal M)\) pätee
  \[
    \mathcal M \models_{s(b/x)(a/y)} A.
  \]
  Tässä sama alkio \(b\) toimii jokaisella \(a\):n arvolla.

  Olkoon nyt \(a \in \dom(\mathcal M)\) mielivaltainen. Koska \(x\) ja
  \(y\) ovat eri muuttujia, tulkintafunktion päivitykset voidaan tehdä
  kummassa tahansa järjestyksessä:
  \[
    s(b/x)(a/y) = s(a/y)(b/x).
  \]
  Siis \(\mathcal M \models_{s(a/y)(b/x)} A\), joten alkio \(b\)
  kelpaa eksistenssikvanttorin todistajaksi ja saadaan
  \[
    \mathcal M \models_{s(a/y)} \exists xA.
  \]
  Koska \(a\) oli mielivaltainen, universaalikvanttorin totuusmääritelmän
  nojalla \(\mathcal M \models_s \forall y\exists xA\).
  Mallin \(\mathcal M\) ja tulkintafunktion \(s\) mielivaltaisuudesta seuraa
  \[
    \exists x\forall yA \models \forall y\exists xA.
  \]
  Todistuksessa ei käytetty mitään oletuksia kaavan \(A\) rakenteesta,
  joten väite pätee kaikilla kaavoilla \(A\).\qedhere
\end{solution}

\begin{exercise}
  Mitkä muuttujien esiintymistä seuraavissa kaavoissa ovat vapaita ja mitkä
  sidottuja? Mitkä kaavoista ovat lauseita?
  \begin{enumerate}[label=(\alph*)]
    \item \(\forall x(P_0(x) \to P_1(y))\)
    \item \(\forall x(\forall yE(x,y) \lor \forall zE(y,z))\)
    \item \(\forall y(\exists x\,R(x,y) \lor \exists x\,R(y,x))\)
    \item \(\forall x(P(x) \to \exists yR(x,y))\)
    \item \(\forall xP_0(y)\)
    \item \(\forall y(\exists x\,x = y \lor R(y,x))\)
  \end{enumerate}
\end{exercise}

\begin{exercise}
  Osoita, että sidotun muuttujan uudelleenkvantifiointi ei tuo kaavaan mitään
  uutta, eli osoita
  \begin{enumerate}[label=(\alph*)]
    \item \(\forall x_i\exists x_iA \Leftrightarrow \exists x_iA\),
    \item \(\exists x_i\forall x_iA \Leftrightarrow \forall x_iA\).
  \end{enumerate}
\end{exercise}

\begin{exercise}
  Jos \(A\) on \(L\)-kaava, niin \(A\):n \emph{universaalisulkeuma} on
  lause, joka saadaan, kun \(A\):n kaikki vapaat muuttujat kvaintifioidaan
  universaalikvanttoreilla. Esim. kaavan \(R(x_0,x_1) \land P(x_0)\)
  universaalisulkeuma on lause
  \(\forall x_0\forall x_1(R(x_0,x_1) \land P(x_0))\).
  \begin{enumerate}[label=(\alph*)]
    \item Osoita, että mielivaltaisessa \(L\)-mallissa \(\mathcal M\),
      \(A\) on tosi jos ja vain jos \(A\):n universaalisulkeuma on tosi.
    \item Anna esimerkki kaavasta \(A\), joka ei ole loogisesti ekvivalentti
      universaalisulkeumansa kanssa. Perustele.
  \end{enumerate}
\end{exercise}

\begin{exercise}
  Todista materiaalin lause 4.7: Jos \(s\) ja \(s'\) ovat kaksi
  \(\mathcal M\)-tulkintafunktiota, joilla \(s(x_i) = s'(x_i)\) kaikilla
  kaavan \(A\) vapailla muuttujilla \(x_i\), niin
  \[
    \mathcal M \models_s A
    \quad\text{jos ja vain jos}\quad
    \mathcal M \models_{s'} A.
  \]
  Voit käyttää todistuksessasi seuraavaa aputulosta, joka on triviaali
  aakkostoille ilman funktiosymboleita, ja jonka voi myöhemmin yleistää
  kaikille aakkostoille:
  \begin{quote}
    Lemma. Jos \(s\) ja \(s'\) ovat kaksi \(\mathcal M\)-tulkintafunktiota,
    joilla \(s(x_i) = s'(x_i)\) kaikilla termissä \(t\) esiintyvillä
    muuttujilla \(x_i\), niin
    \(t^{\mathcal M}\langle s\rangle = t^{\mathcal M}\langle s'\rangle\).
  \end{quote}
\end{exercise}

\begin{exercise}
  Minkä joukon kaava
  \(\exists x_1\exists x_2((E(x_0,x_1) \land E(x_1,x_2)) \land \neg x_0 = x_2)\)
  määrittelee allaolevassa verkossa?
  \begin{center}
    \includegraphics[width=0.4\linewidth]{figures/exercise-03-graph.pdf}
  \end{center}
\end{exercise}

\begin{exercise}
  Alkio \(a\) on määriteltävä mallissa \(\mathcal M\) jos joukko \(\{a\}\)
  on määriteltävä \(\mathcal M\):ssä. Osoita, että \(2\) on määriteltävä
  mallissa \((\NN,<)\), missä \(<\) on luonnollisten lukujen luonnollinen
  järjestys.
\end{exercise}

\end{document}
